\documentclass[12pt, a4paper, oneside]{report}

% --- Packages de base ---
\usepackage[utf8]{inputenc}
\usepackage[T1]{fontenc}
\usepackage[french]{babel}
\usepackage{libertine} % Police plus professionnelle
\usepackage[libertine]{newtxmath}
\usepackage{microtype} % Améliore la justification du texte
\usepackage{enumitem}
% --- Mise en page et Géométrie ---
\usepackage{geometry}
\usepackage{float}
\usepackage[utf8]{inputenc}
\usepackage[T1]{fontenc}
\usepackage{array}
\usepackage{float}
\geometry{
    a4paper,
    total={160mm,245mm},
    left=30mm,   % Marge plus large à gauche pour la reliure
    top=25mm,
  \begin{figure}[h]
    \centering
    \includegraphics[width=0.8\linewidth]{cdr7.png}
    \caption[Exemple des doublons détectés]{Exemple des doublons détectés. Vue de l'interface de gestion des doublons CDR montrant une table récapitulative des groupes d'enregistrements répétés (colonnes : source, MSISDN, service, date & heure, montant, occurrences, impact, actions). La colonne « occurrences » indique le nombre de répétitions (ex : 3x, 5x) et permet de classer la criticité. Cette figure illustre également des éléments d'aide à la décision : identification des services les plus affectés, visualisation des périodes à risque et possibilités d'action (purge unitaire, purge en masse, export pour enquête).}
    \label{fig:duplicates-example}
  \end{figure}
\definecolor{mainblue}{RGB}{0, 68, 148} % Bleu Tunisie Telecom
\usepackage[colorlinks=true, linkcolor=black, citecolor=mainblue, urlcolor=mainblue]{hyperref}

% --- Graphiques et Tableaux ---
\usepackage{graphicx}
\usepackage{float}
\usepackage{booktabs} % Pour des tableaux style professionnel
\usepackage{array}
\usepackage{longtable}
\usepackage{caption}
\captionsetup{labelfont={bf,color=mainblue}, font=small}

% --- En-têtes et Pieds de page ---
\usepackage{fancyhdr}
\pagestyle{fancy}
\fancyhf{}
\fancyhead[L]{\small \leftmark}
\fancyhead[R]{\small \thepage}
\renewcommand{\headrulewidth}{0.5pt}

% --- Style des Titres ---
\usepackage{titlesec}
\titleformat{\chapter}[display]
  {\normalfont\huge\bfseries\color{mainblue}}
  {\chaptertitlename\ \thechapter}{20pt}{\Huge}
\titlespacing*{\chapter}{0pt}{-20pt}{40pt}

\titleformat{\section}
  {\normalfont\Large\bfseries\color{mainblue}}
  {\thesection}{1em}{}

% --- Espacement ---
\usepackage{setspace}
\onehalfspacing % Interligne 1.5 obligatoire pour les rapports

% --- Code source ---
\usepackage{listings}
\lstset{
    basicstyle=\ttfamily\small,
    keywordstyle=\color{blue},
    commentstyle=\color{gray},
    breaklines=true,
    frame=single,
    backgroundcolor=\color{gray!5}
}

% [Contenu du document similaire...]
% Pour brevité, j'ai omis les parties identiques

% --- Sprint 3: Réconciliation et Visualisation des données ---
\chapter{Sprint 3 : Réconciliation et Visualisation des données}

\section{Introduction}
Ce premier sprint est dédié à la phase d'immersion technique et fonctionnelle...

% [Reste du contenu du Sprint 3...]

\subsection{Conclusion du sprint}
Le Sprint 3 a marqué une étape décisive dans le projet en transformant des volumes de données brutes en une intelligence métier exploitable et visuelle. L'implémentation des algorithmes de réconciliation a permis de lever le voile sur les disparités entre les flux MMG et OCC, instaurant ainsi un climat de confiance technique grâce à la détection automatisée des écarts de trafic.

Sur le plan technologique, le déploiement de l'interface sous React a apporté une dimension interactive indispensable à l'expérience utilisateur. L'adoption d'une architecture modulaire et l'intégration de Recharts ont permis de traduire des KPIs complexes en visualisations intuitives, facilitant ainsi la prise de décision rapide pour les gestionnaires.

Enfin, la phase rigoureuse de tests et de validation a confirmé la robustesse de la solution. En confrontant les composants front-end et les services back-end à des scénarios de charge réels, nous avons pu stabiliser le système. Ce sprint s'achève sur une plateforme fonctionnelle, prête à offrir une transparence totale sur les revenus et les anomalies opérationnelles.

\newpage

% --- NOUVEAU CHAPITRE 5 : CONCLUSION GÉNÉRALE ---
\chapter{Conclusion générale}

Arrivé au terme de ce projet de fin d'études mené à la Direction Centrale des Finances de Tunisie Telecom, nous avons créé une plateforme web avancée pour la gestion électronique et la réconciliation des courriers et des flux financiers SMS+ (VAS).

Ce projet répondait à des défis majeurs : centraliser des données hétérogènes, éradiquer les anomalies de facturation multi-opérateurs, et moderniser les outils d'audit financier.

Nous avons adopté une démarche de développement agile pour structurer notre progression de manière incrémentale.

Les premières phases ont permis de créer un pipeline ETL conteneurisé sous Docker pour ingérer, nettoyer et structurer des volumes industriels de tickets d'appels (CDR) dans une base de données PostgreSQL.

Ensuite, nous avons implémenté des processus de déduplication et des moteurs de réconciliation quantitative et qualitative pour isoler en temps réel les écarts de trafic.

La plateforme dispose désormais d'une interface utilisateur réactive et modulaire sous React et Recharts, couplée à une couche d'Intelligence Artificielle hybride.

Cette expérience a été très enrichissante, tant sur le plan personnel que professionnel. Nous avons consolidé nos compétences en ingénierie logicielle full-stack et appris à collaborer avec des experts métiers.

Même si la plateforme est pleinement fonctionnelle, nous envisageons des perspectives d'évolution, notamment le traitement en streaming temps réel, l'apprentissage automatique local approfondi, et l'extension fonctionnelle.

Ce projet montre qu'en combinant des architectures web modernes et l'intelligence artificielle, il est possible de transformer des volumes massifs de données en un levier d'audit financier transparent et robuste.

\newpage

% --- Annexe ---
\appendix
\chapter{Annexe}
\section{Création d'une VM}
Ces figures présentent les étapes de création de la machine virtuelle.

\section{Installation et configuration Hadoop}
Ces figures présentent les étapes de l'installation et la configuration de Hadoop.

\section{Installation et configuration de Apache Nifi}
Ces figures présentent les étapes de l'installation et la configuration de Apache Nifi.

\section{Installation et configuration de Apache Spark}
Ces figures présentent les étapes de l'installation et la configuration de Apache Spark.

\section{Installation et configuration de MongoDB}
Ces figures présentent les étapes de l'installation et la configuration de MongoDB.

\section{Configuration de Power BI avec la VM}
Ces figures présentent les étapes de la configuration de Power BI.

% --- Webographie ---
\chapter{Webographie}
\begin{thebibliography}{9}
    \bibitem{tt}
    Tunisie Télécom: Historique, services fixes et mobiles, etc. (it-tunisie.tn).

    \bibitem{org_tt}
    \url{https://www.researchgate.net/figure/Organization-chart-of-Tunisie-Telecom1_fig1_345952324}.

    \bibitem{scrum}
    \url{https://www.planzone.fr/blog/quest-ce-que-la-methodologie-agile}.

    \bibitem{bigdata}
    \url{https://www.oracle.com/fr/big-data/what-is-big-data/}.

    \bibitem{hadoop}
    \url{https://ledatascientist.com/hadoop-pour-les-nuls-presentation-de-lecosysteme-hadoop/}.

\end{thebibliography}

\end{document}
